
\begin{proposition}[Normalization, conductor, and the excess kernel]\label{cdp:st:A.2.2}
Let $\eta:\widetilde S\to S$ be the finite analytic normalization, $\mathcal B=\eta_*\OO_{\widetilde S}$, and $\mathfrak c=\Ann_{\OO_S}(\mathcal B/\OO_S)$ its conductor. There are canonical inclusions
\[
 \mathcal J\subset\mathfrak c,\qquad
 \mathcal B\subset\mathcal J^{-1},\qquad
 \eta_*\eta^*N^*\hookrightarrow K.
\tag{A.3}\label{cdp:eq:A.3}
\]
In a local ring $R$ with normalization $B$, put $D=\Hom_B(JB,B)$, viewed as a fractional ideal. Then
\[
 J^{-1}=\ker\left[D\longrightarrow
 \Hom_R\bigl(J,(B/\mathfrak c)/(R/\mathfrak c)\bigr)\right],
\quad q\longmapsto\bigl(b\mapsto qb\bmod R\bigr).
\tag{A.4}\label{cdp:eq:A.4}
\]
For any locally free $\OO_S$-module $E$, the conductor gives a second, independent exact section criterion:
\begin{multline}
H^0(S,E)=
\ker\left[H^0(\widetilde S,\eta^*E)\oplus H^0(D_c,E|_{D_c})\right.\\
\left.\longrightarrow H^0(\widetilde D_c,\eta^*E|_{\widetilde D_c})\right],
\qquad (a,b)\longmapsto a|_{\widetilde D_c}-\eta^*b,
\tag{A.5}\label{cdp:eq:A.5}
\end{multline}
where $D_c$ and $\widetilde D_c$ are defined by $\mathfrak c$ in $S$ and $\widetilde S$. In particular, the ambient group vanishes exactly when the map in \eqref{cdp:eq:A.5}, for $E=V\otimes A$, is injective.
\end{proposition}
\begin{proof}
We prove $JB\subset R$ rather than assume a conductor formula. Choose a generic linear projection for which Weierstrass preparation writes the equation as a unit times a monic polynomial $P(z)$ over $A_0=\C\{x_1,\ldots,x_n\}$. Thus $R=A_0[z]/(P)$ is finite free over $A_0$ and generically separable. The functional taking the coefficient of $z^{m-1}$ in a remainder modulo $P$, where $m=\deg P$, generates $\Hom_{A_0}(R,A_0)$ as an $R$-module: the matrix of the pairing on $1,z,\ldots,z^{m-1}$ has zeros below the reverse diagonal and ones on that diagonal. Over the fraction field, Lagrange interpolation identifies that functional with $r\mapsto\operatorname{Tr}(r/P'(z))$. Therefore the trace dual is $P'(z)^{-1}R$.

For $b\in B$ and $r\in R$, the element $br$ is integral over the integrally closed ring $A_0$. Its trace belongs to $A_0$. It follows that $b$ lies in this trace dual, so $P'(z)b\in R$. Modulo $F$, the derivative in the chosen projection direction differs from $P'(z)$ by a unit. The permissible projection directions form a nonempty open set; choosing a linear basis of such directions proves $F_{x_i}B\subset R$ for every $i$. This proves $J\subset\mathfrak c$ and hence $B\subset J^{-1}$. The last inclusion of \eqref{cdp:eq:A.3} is (A.2) and the projection formula.

The module $D$ consists exactly of the meromorphic $q$ satisfying $qJ\subset B$. Requiring that these products descend to $R$ is precisely the vanishing of their classes in $B/R=(B/\mathfrak c)/(R/\mathfrak c)$, proving \eqref{cdp:eq:A.4}. Finally the ring map
\[
 R\longrightarrow B\oplus R/\mathfrak c,
 \qquad r\longmapsto(r,r\bmod\mathfrak c)
\]
has image the kernel of $(b,\bar r)\mapsto b-\bar r\in B/\mathfrak c$. Surjectivity onto $B/\mathfrak c$ comes from the first summand. Tensoring this exact sequence by the free module corresponding to $E$ and taking global sections gives \eqref{cdp:eq:A.5}. All conductor restrictions are scheme-theoretic, so this argument includes nonreduced conductors.
\end{proof}

The fractional ideal appearing in (A.2) has an established interpretation. Granger--Schulze, Proposition~3.4 and Corollary~3.5 of \cite{GS2014}, identify $J^{-1}$ with the module of logarithmic residues and prove the inclusions in \eqref{cdp:eq:A.3}. The preceding proof gives the maps needed here directly. Their Theorems~1.1 and~1.3 also give the following exact extension of the normal-crossings case.


